\chapter{Introduction}
\enlargethispage{2\baselineskip}

\section{Objective and scope}

The application generates new speech in the voice of a reference speaker on
Android. It takes three inputs: a reference recording, a transcript containing
the words spoken in that recording, and the target text to synthesize. After
installation, generation runs locally on the phone. The resulting audio can be
played, revisited in the app's library, exported as a WAV file or shared through
Android.

The main engineering constraints are the size of the two neural models, the
memory and compute available on a phone, and the fixed audio window that must
fit both the reference and the generated speech. The project adapts and deploys
pretrained IndicF5 and Vocos. The contribution is the conversion and native
deployment pipeline, including reference audio processing, multilingual text
segmentation and joining generated waveforms for longer passages.

\section{Earlier project stages}

\begin{enumerate}[nosep]
  \item \textbf{English Android baseline.} The Flutter prototype used ZipVoice-Distill
        through Sherpa-ONNX. It established that a phone could record or import a
        reference clip, synthesize new speech locally and play the result.
        \texttt{tts\_service.dart} ran the English inference work in a Dart isolate.
  \item \textbf{IndicF5 laptop proof of concept.} The PyTorch trials showed that
        the pretrained model could synthesize examples in Indian languages.
        They also identified the need to budget target duration against the
        reference. \texttt{clone.py} exercised individual prompts and
        \texttt{batch.py} ran multiple prompts.
  \item \textbf{Mobile model export.} Static INT8 and dynamic FP32 ONNX
        variants were investigated. The current Android deployment instead
        converts the PyTorch models directly to LiteRT and manages inference
        through Kotlin and the native C++ engine.
\end{enumerate}
